\begin{example}
    La sucesión $x=(x_n)_{n\geq1}$ dada por
    \[
        0, \ 
        0, \ \frac{1}{2}, \ 
        0, \ \frac{1}{3}, \ \frac{2}{3}, \ 
        0, \ \frac{1}{4}, \ \frac{2}{4}, \ \frac{3}{4}, \ 
        0, \ \frac{1}{5}, \ \frac{2}{5}, \ \frac{3}{5}, \ \frac{4}{5}, \ 
        0, \ \dots 
    \]
    es equidistribuida módulo $1$.
\end{example}

\begin{resolve}
    Para cada $k \in \mathbb{N}$, consideramos los conjuntos $A_k = \left\{\frac{i}{k} : i = 0, 1, 2, \dots, k-1 \right\}$.
    La sucesión se puede escribir como la concatenación de estos conjuntos,
    \[
        A_1, \ A_2, \ A_3, \ A_4, \ A_5, \ \dots
    \]
    Sea $(a,b)$ con $a,b \notin \mathbb{Q}$, sea $k$ fijo, estudiamos cuántos elementos de $A_k$ pertenecen al intervalo,
    \[
        P_k = 
        \card \left\{ \frac{i}{k} \in A_k : \frac{i}{k} \in (a,b) \right\} =
        \card \left\{ \frac{i}{k} \in A_k : i \in (ka,kb)\right\}
    \]
    Contando el número de enteros en el intervalo $(ka,kb)$, obtenemos que 
    \[
        P_k = \floor{kb} - \floor{ka+1} + 1 = \floor{kb} - \floor{ka}
    \]

    Tomando la sucesión hasta el conjunto $A_M$, inclusive. El número de total de términos y la cantidad
    de términos que están en el intervalo $(a,b)$ es
    \[
        T_M = \sum_{k=1}^M k = \frac{M(M+1)}{2} \qquad 
        S_M = \sum_{k=1}^M P_k = \sum_{k=1}^M \left(\floor{kb} - \floor{ka}\right)
    \]

    Veamos que, contando por conjuntos, se cumple la definición de equidistribución, es decir,
    \[
        \left|\frac{S_M}{T_M} - (b-a)\right| \longrightarrow 0 \qquad \text{cuando } M \to \infty.
    \]

    Para ello, tratamos de acotarlo usando que $\partfrac{x} \in [0,1)$ para cada $x \in \mathbb{R}$,
    \begin{align*}
        \left| \frac{S_M}{T_M} - (b-a) \right| 
        =\left| \frac{S_M - T_M(b-a)}{T_M} \right| 
        &=\frac{1}{T_M}\left| \sum_{k=1}^M \left(\floor{kb} - \floor{ka}\right) - \sum_{k=1}^M k(b-a) \right| \\ 
        &=\frac{1}{T_M}\left| \sum_{k=1}^M \left(\floor{kb} - \floor{ka} - kb + ka\right) \right| \\
        &\leq\frac{1}{T_M} \sum_{k=1}^M \left|\partfrac{ka} - \partfrac{kb}\right| 
        \leq\frac{1}{T_M} \sum_{k=1}^M 1 = \frac{M}{T_M}
    \end{align*}
    
    Por lo tanto, 
    \[
        \left| \frac{S_M}{T_M} - (b-a) \right| \leq
        \frac{M}{T_M} = \frac{2}{M+1} \longrightarrow 0 \qquad \text{cuando } M \to \infty.
    \]
    
    Esto demuestra que, al contar por conjuntos, la sucesión es equidistribuida para $N = T_M$.
    Falta~comprobar que también lo es para cualquier $N$. 

    Sea $N$ tal que $T_M \leq N \leq T_{M+1}$,
    \[
        A_N((a,b), x) = S_M + R \qquad \text{con } 0 \leq R \leq P_{M+1}
    \]
    Tenemos las siguientes desigualdades,
    \[
        S_M \leq S_M + R \leq S_{M+1} 
        \quad \text{ y } \quad
        \frac{1}{T_{M+1}} \leq \frac{1}{N} \leq \frac{1}{T_M}
    \]
    Al multiplicarlas,
    \[
        \frac{S_M}{T_{M+1}} \leq \frac{S_M + R}{N} \leq \frac{S_{M+1}}{T_M}
    \]
    
    Al tomar límites cuando $M \to \infty$, y por lo tanto $N \to \infty$, obtenemos
    \begin{align*}
        \lim_{M \to \infty} \frac{T_M}{T_{M+1}} \, \frac{S_M}{T_M} \leq
        \liminf_{N \to \infty} \frac{S_M + R}{N} &\leq
        \limsup_{N \to \infty} \frac{S_M + R}{N} \leq
        \lim_{M \to \infty} \frac{T_{M+1}}{T_M} \, \frac{S_{M+1}}{T_{M+1}}  \\
        b-a \leq
        \liminf_{N \to \infty} \frac{S_M + R}{N} &\leq
        \limsup_{N \to \infty} \frac{S_M + R}{N} \leq
        b-a
    \end{align*}
    Por lo que, el límite existe y es
    \[
        \lim_{N \to \infty} \frac{A_N((a,b), x)}{N} =
        \lim_{N \to \infty} \frac{S_M + R}{N} = 
        b-a
    \]
\end{resolve}